\documentclass[11pt,a4paper]{article}

% ── 한글 지원 ──
\usepackage{kotex}
\usepackage{fontspec}

% ── 레이아웃 ──
\usepackage[margin=2.5cm]{geometry}
\usepackage{setspace}
\onehalfspacing

% ── 패키지 ──
\usepackage{booktabs}
\usepackage{longtable}
\usepackage{array}
\usepackage{tabularx}
\usepackage{xcolor}
\usepackage{listings}
\usepackage{hyperref}
\usepackage{enumitem}
\usepackage{fancyhdr}
\usepackage{titlesec}
\usepackage{tocloft}
\usepackage{caption}

% ── 색상 정의 ──
\definecolor{codebg}{RGB}{245,245,245}
\definecolor{codeframe}{RGB}{200,200,200}
\definecolor{keyword}{RGB}{0,100,170}
\definecolor{string}{RGB}{170,55,0}
\definecolor{comment}{RGB}{100,130,100}
\definecolor{accent}{RGB}{0,80,150}
\definecolor{lightaccent}{RGB}{230,240,250}

% ── 코드 블록 설정 ──
\lstset{
    language=Python,
    basicstyle=\small\ttfamily,
    keywordstyle=\color{keyword}\bfseries,
    stringstyle=\color{string},
    commentstyle=\color{comment}\itshape,
    backgroundcolor=\color{codebg},
    frame=single,
    rulecolor=\color{codeframe},
    breaklines=true,
    breakatwhitespace=true,
    showstringspaces=false,
    tabsize=4,
    xleftmargin=1em,
    xrightmargin=1em,
    aboveskip=1em,
    belowskip=1em,
    literate={→}{$\rightarrow$}1
             {─}{-}1
             {│}{|}1,
}

\lstdefinestyle{bash}{
    language=bash,
    basicstyle=\small\ttfamily,
    backgroundcolor=\color{codebg},
    frame=single,
    rulecolor=\color{codeframe},
    xleftmargin=1em,
    xrightmargin=1em,
    aboveskip=1em,
    belowskip=1em,
    breaklines=true,
}

% ── 하이퍼링크 설정 ──
\hypersetup{
    colorlinks=true,
    linkcolor=accent,
    urlcolor=accent,
    citecolor=accent,
    bookmarks=true,
    bookmarksopen=true,
}

% ── 헤더/푸터 ──
\pagestyle{fancy}
\fancyhf{}
\fancyhead[L]{\small\textcolor{gray}{글로벌 거시경제 데이터 수집 시스템}}
\fancyhead[R]{\small\textcolor{gray}{\thepage}}
\renewcommand{\headrulewidth}{0.4pt}

% ── 섹션 스타일 ──
\titleformat{\section}
    {\Large\bfseries\color{accent}}
    {\thesection}{1em}{}[\vspace{-0.5em}\textcolor{accent}{\rule{\textwidth}{0.5pt}}]

\titleformat{\subsection}
    {\large\bfseries\color{accent!80!black}}
    {\thesubsection}{1em}{}

\titleformat{\subsubsection}
    {\normalsize\bfseries\color{accent!60!black}}
    {\thesubsubsection}{1em}{}

% ── 새 명령어 ──
\newcommand{\apibox}[1]{%
    \noindent\colorbox{lightaccent}{\parbox{\dimexpr\textwidth-2\fboxsep}{%
        \small\texttt{#1}%
    }}%
    \vspace{0.5em}%
}

\newcommand{\param}[1]{\texttt{\textcolor{keyword}{#1}}}
\newcommand{\code}[1]{\texttt{#1}}
\newcommand{\file}[1]{\texttt{\textcolor{string}{#1}}}

% ══════════════════════════════════════════════════════════════
\begin{document}

% ── 표지 ──
\begin{titlepage}
\centering
\vspace*{4cm}

{\Huge\bfseries\textcolor{accent}{글로벌 거시경제 데이터}\par}
\vspace{0.5cm}
{\Huge\bfseries\textcolor{accent}{수집 시스템 매뉴얼}\par}

\vspace{2cm}

{\Large FRED $\cdot$ OECD $\cdot$ BIS $\cdot$ IMF $\cdot$ World Bank $\cdot$ CIA Factbook $\cdot$ ECOS\par}

\vspace{3cm}

{\large \code{download\_global.py} 사용 가이드\par}

\vspace{2cm}

{\large 김남일\par}
{\normalsize DS증권 퀀트\par}
\vspace{0.3cm}
{\large NAMIL KIM\par}
{\normalsize DS Investment \& Securities\par}

\vspace{1.5cm}

{\large 2026년 2월\par}

\vfill

{\normalsize
저장 경로: \code{data/}\\[0.3em]
출력 형식: Apache Parquet (.parquet) + CSV (.csv)
}
\end{titlepage}

% ── 목차 ──
\newpage
\tableofcontents
\newpage


% ══════════════════════════════════════════════════════════════
\section{개요}
% ══════════════════════════════════════════════════════════════

이 시스템은 6개 국제기관 API + 한국은행 ECOS API를 통해 글로벌 거시경제 데이터를
자동으로 수집하고 Apache Parquet 형식으로 저장합니다.
모든 Python 모듈(\code{*.py})과 \code{data/} 디렉토리가 동일한 작업 디렉토리에 위치합니다.

\subsection{시스템 구성}

\begin{table}[h]
\centering
\caption{API 모듈 구성}
\begin{tabularx}{\textwidth}{l l X l}
\toprule
\textbf{소스} & \textbf{모듈 파일} & \textbf{데이터 범위} & \textbf{인증} \\
\midrule
FRED     & \file{fred\_api.py}      & 미국 GDP, CPI, 고용, 금리 등      & API 키 필요 \\
OECD     & \file{oecd\_api.py}      & 회원국 GDP, CPI, 실업률, 금리      & 불필요 \\
BIS      & \file{bis\_api.py}       & CPI, 정책금리, 신용, 부동산, 환율   & 불필요 \\
IMF      & \file{imf\_api.py}       & WEO: GDP, 인플레이션, 재정, 경상수지 & 불필요 \\
World Bank & \file{worldbank\_api.py} & GDP, CPI, 실업률, 무역, 인구 등    & 불필요 \\
Factbook & \file{factbook\_api.py}  & 국가별 경제 스냅샷 (최신)           & 불필요 \\
ECOS     & \file{ecos\_api.py}      & 한국은행 경제통계 (163개 테이블)     & API 키 내장 \\
\bottomrule
\end{tabularx}
\end{table}

\subsection{저장 형식}

\begin{itemize}[leftmargin=2em]
    \item \textbf{경로}: \code{data/} (스크립트와 동일 디렉토리 하위)
    \item \textbf{기본 형식}: Apache Parquet (컬럼 기반 압축 포맷)
    \item \textbf{CSV 형식}: \code{--csv} 옵션 사용 시 Parquet과 함께 CSV도 동시 저장 (UTF-8 BOM)
    \item \textbf{파일명}: \code{\{소스\}\_\{지표\}.parquet} / \code{\{소스\}\_\{지표\}.csv}
    \item \textbf{기간}: 기본 2007년~현재 (글로벌 금융위기 이후)
\end{itemize}

\subsection{대상 국가}

\begin{table}[h]
\centering
\caption{소스별 대상 국가}
\begin{tabularx}{\textwidth}{l l X}
\toprule
\textbf{소스} & \textbf{코드 체계} & \textbf{기본 대상국} \\
\midrule
FRED     & 미국 전용   & 미국(US) \\
OECD     & ISO 3자리   & KOR, USA, JPN, DEU, GBR, FRA, CHN, CAN, AUS, ITA \\
BIS      & ISO 2자리   & KR, US, JP, DE, GB, FR, CN, CA, AU, IT \\
IMF      & ISO 3자리   & USA, KOR, JPN, DEU, GBR, FRA, CHN, ITA, CAN, AUS, MEX, BRA, IND, IDN, TUR \\
World Bank & ISO 2자리 & US, KR, JP, DE, GB, FR, CN, IT, CA, AU, MX, BR, IN, ID, TR \\
Factbook & GEC(FIPS)   & us, ks, ja, gm, uk, fr, ch, it, ca, as, mx, br, in, id, tu \\
\bottomrule
\end{tabularx}
\end{table}


% ══════════════════════════════════════════════════════════════
\section{빠른 시작}
% ══════════════════════════════════════════════════════════════

\subsection{사전 요구사항}

\begin{lstlisting}[style=bash]
# Python 패키지
pip install requests pandas pyarrow

# FRED API 키 설정 (FRED 데이터 필요시)
export FRED_API_KEY=your_api_key_here
# 무료 발급: https://fred.stlouisfed.org/docs/api/api_key.html
\end{lstlisting}

\subsection{전체 다운로드}

\begin{lstlisting}[style=bash]
# 모든 소스 다운로드 (parquet만)
python3 download_global.py

# parquet + csv 동시 저장
python3 download_global.py --csv

# 선택적 다운로드
python3 download_global.py fred oecd        # FRED와 OECD만
python3 download_global.py --csv imf wb     # IMF와 WB (parquet + csv)
python3 download_global.py --csv bis        # BIS (parquet + csv)
\end{lstlisting}

\subsection{소스 선택 옵션}

\begin{table}[h]
\centering
\caption{\code{download\_global.py} 실행 인자}
\begin{tabular}{l l l}
\toprule
\textbf{인자} & \textbf{구분} & \textbf{설명} \\
\midrule
\multicolumn{3}{l}{\textbf{소스 선택} (생략 시 전체)} \\
\midrule
\code{fred}     & 소스 & FRED (미국 연방준비) \\
\code{oecd}     & 소스 & OECD (경제협력개발기구) \\
\code{bis}      & 소스 & BIS (국제결제은행) \\
\code{imf}      & 소스 & IMF (국제통화기금) \\
\code{wb}       & 소스 & World Bank (세계은행) \\
\code{factbook} & 소스 & CIA Factbook \\
\midrule
\multicolumn{3}{l}{\textbf{옵션}} \\
\midrule
\code{-{}-csv}  & 플래그 & Parquet과 함께 CSV도 동시 저장 (UTF-8 BOM) \\
\bottomrule
\end{tabular}
\end{table}

인자 없이 실행하면 전체 6개 소스를 순차적으로 다운로드합니다 (Parquet만).
\code{--csv} 플래그를 추가하면 각 항목마다 \file{.parquet}과 \file{.csv}를 모두 저장합니다.
FRED API 키가 없으면 FRED는 자동 건너뜁니다.

\subsection{결과 확인}

다운로드 완료 후 결과 요약이 출력됩니다:

\begin{lstlisting}[style=bash]
# 저장된 파일 확인
ls data/*.parquet | wc -l    # parquet 파일 수
ls data/*.csv | wc -l        # csv 파일 수 (--csv 사용 시)

# Python에서 로드 (parquet)
python3 -c "
import pandas as pd
df = pd.read_parquet('data/fred_gdp_real.parquet')
print(df.tail())
"

# Python에서 로드 (csv)
python3 -c "
import pandas as pd
df = pd.read_csv('data/fred_gdp_real.csv')
print(df.tail())
"
\end{lstlisting}


% ══════════════════════════════════════════════════════════════
\section{FRED API (\file{fred\_api.py})}
% ══════════════════════════════════════════════════════════════

\subsection{개요}

\textbf{FRED} (Federal Reserve Economic Data)는 미국 세인트루이스 연방준비은행이 운영하는
거시경제 데이터베이스로, 미국 중심의 816,000+개 시계열을 제공합니다.

\begin{itemize}[leftmargin=2em]
    \item \textbf{공식 사이트}: \url{https://fred.stlouisfed.org/}
    \item \textbf{API 문서}: \url{https://fred.stlouisfed.org/docs/api/fred/}
    \item \textbf{API 키}: \url{https://fred.stlouisfed.org/docs/api/api_key.html}에서 무료 발급
    \item \textbf{기본 URL}: \code{https://api.stlouisfed.org/fred}
\end{itemize}

\subsection{API 키 설정}

\begin{lstlisting}[language=Python]
# 방법 1: 환경변수
export FRED_API_KEY=your_key_here

# 방법 2: 코드에서 직접
from fred_api import set_api_key
set_api_key("your_key_here")
\end{lstlisting}

\subsection{핵심 함수}

\subsubsection{\code{get\_series()} --- 단일 시계열 조회}

\begin{lstlisting}[language=Python]
get_series(series_id, start_date=None, end_date=None,
           frequency=None, units=None, sort_order="asc")
\end{lstlisting}

\begin{table}[h]
\centering
\begin{tabularx}{\textwidth}{l l X}
\toprule
\textbf{매개변수} & \textbf{타입} & \textbf{설명} \\
\midrule
\param{series\_id} & str & 시리즈 ID (예: \code{"GDPC1"}, \code{"CPIAUCSL"}) \\
\param{start\_date} & str & 시작일 \code{YYYY-MM-DD} \\
\param{end\_date} & str & 종료일 \code{YYYY-MM-DD} \\
\param{frequency} & str & 주기 변환: \code{d/w/bw/m/q/sa/a} \\
\param{units} & str & 단위 변환 (아래 표 참조) \\
\bottomrule
\end{tabularx}
\end{table}

\textbf{units 옵션:}
\begin{table}[h]
\centering
\begin{tabular}{l l}
\toprule
\textbf{코드} & \textbf{의미} \\
\midrule
\code{lin} & 원본 (기본값) \\
\code{chg} & 전기 변화분 \\
\code{pch} & 전기 변화율 (\%) \\
\code{pc1} & 전년 동기 변화율 (\%) \\
\code{pca} & 연율화 변화율 (\%) \\
\code{log} & 자연로그 \\
\bottomrule
\end{tabular}
\end{table}

\textbf{반환}: \code{DataFrame [date, value, series\_id]}

\subsubsection{\code{get\_multi\_series()} / \code{get\_multi\_series\_wide()}}

\begin{lstlisting}[language=Python]
# Long format (date, value, series_id)
df = get_multi_series(["GDPC1", "CPIAUCSL", "UNRATE"], "2020-01-01")

# Wide format (date index, series_id columns)
wide = get_multi_series_wide(["GDPC1", "CPIAUCSL", "UNRATE"], "2020-01-01")
\end{lstlisting}

\subsection{다운로드 항목 (17개)}

\begin{longtable}{l l l}
\toprule
\textbf{파일명} & \textbf{시리즈 ID} & \textbf{설명} \\
\midrule
\endhead
\file{fred\_gdp\_real}         & GDPC1           & 실질 GDP (분기, 연율) \\
\file{fred\_gdp\_growth}       & A191RL1Q225SBEA & 실질 GDP 성장률 (\%) \\
\file{fred\_cpi}               & CPIAUCSL        & CPI 전체 (월별, 계절조정) \\
\file{fred\_cpi\_core}         & CPILFESL        & Core CPI (식품·에너지 제외) \\
\file{fred\_pce}               & PCEPI           & PCE 물가지수 (월별) \\
\file{fred\_pce\_core}         & PCEPILFE        & Core PCE (월별) \\
\file{fred\_unemployment}      & UNRATE          & 실업률 (월별, \%) \\
\file{fred\_payrolls}          & PAYEMS          & 비농업 고용자 수 (천명) \\
\file{fred\_participation}     & CIVPART         & 경제활동참가율 (\%) \\
\file{fred\_initial\_claims}   & ICSA            & 신규 실업수당 청구 (주별) \\
\file{fred\_fed\_funds}        & FEDFUNDS        & 연방기금금리 (월별) \\
\file{fred\_treasury\_10y}     & DGS10           & 10년물 국채수익률 (월별) \\
\file{fred\_treasury\_2y}      & DGS2            & 2년물 국채수익률 (월별) \\
\file{fred\_term\_spread}      & T10Y2Y          & 장단기 스프레드 10Y-2Y \\
\file{fred\_industrial\_prod}  & INDPRO          & 산업생산지수 (월별) \\
\file{fred\_retail\_sales}     & RSXFS           & 소매판매 (월별) \\
\file{fred\_consumer\_sentiment} & UMCSENT       & 소비자심리지수 (월별) \\
\file{fred\_m2}                & M2SL            & M2 통화량 (월별) \\
\bottomrule
\end{longtable}

\subsection{추가 시리즈 사전 (\code{SERIES} dict)}

\file{fred\_api.py}에 40+개 주요 시리즈 ID가 사전으로 정의되어 있습니다:

\begin{lstlisting}[language=Python]
from fred_api import SERIES
print(SERIES["GDP_real"])       # "GDPC1"
print(SERIES["fed_funds_rate"]) # "FEDFUNDS"
print(SERIES["sp500"])          # "SP500"
print(SERIES["vix"])            # "VIXCLS"
\end{lstlisting}


% ══════════════════════════════════════════════════════════════
\section{OECD API (\file{oecd\_api.py})}
% ══════════════════════════════════════════════════════════════

\subsection{개요}

\textbf{OECD SDMX REST API}는 OECD 회원국 및 주요 비회원국의 거시경제 패널 데이터를 제공합니다.
인증 불필요하나, Rate Limit (분당 10회 정도)에 주의해야 합니다.

\begin{itemize}[leftmargin=2em]
    \item \textbf{기본 URL}: \code{https://sdmx.oecd.org/public/rest}
    \item \textbf{응답 형식}: SDMX CSV 2.0
    \item \textbf{국가 코드}: ISO 3자리 (KOR, USA, JPN, DEU, GBR, FRA, CHN, ...)
\end{itemize}

\subsection{데이터플로우 차원 구조}

OECD SDMX API는 데이터플로우마다 차원(dimension)이 다릅니다:

\subsubsection{QNA --- GDP 성장률 (13차원)}

\apibox{FREQ.ADJUSTMENT.REF\_AREA.SECTOR.COUNTERPART\_SECTOR.TRANSACTION.\\
INSTR\_ASSET.ACTIVITY.EXPENDITURE.UNIT\_MEASURE.PRICE\_BASE.TRANSFORMATION.TABLE\_IDENTIFIER}

\textbf{GDP 성장률 예시 키}:
\begin{lstlisting}
Q.Y.KOR+USA+JPN.S1.S1.B1GQ._Z._Z._Z.PC.L.GY.T0102
\end{lstlisting}

여기서 \code{GY}=전년동기비, \code{G1}=전기비, \code{GCM}=전기비 연율.

\subsubsection{KEI --- CPI, 실업률, 금리 등 (7차원)}

\apibox{REF\_AREA.FREQ.MEASURE.UNIT\_MEASURE.ACTIVITY.ADJUSTMENT.TRANSFORMATION}

\begin{table}[h]
\centering
\caption{KEI 주요 MEASURE 코드}
\begin{tabular}{l l}
\toprule
\textbf{코드} & \textbf{설명} \\
\midrule
\code{CP}     & 소비자물가지수 (CPI) \\
\code{UNEMP}  & 실업률 \\
\code{EMP}    & 고용 \\
\code{PRVM}   & 산업생산 (제조업) \\
\code{RS}     & 소매판매 \\
\code{EX} / \code{IM} & 수출 / 수입 \\
\code{IR3TIB} & 3개월 인터뱅크 금리 \\
\code{IRLT}   & 장기 국채수익률 \\
\code{IRSTCI} & 단기 정책금리 \\
\bottomrule
\end{tabular}
\end{table}

\subsubsection{CLI --- 경기선행지수 (5차원)}

\apibox{REF\_AREA.FREQ.MEASURE.UNIT\_MEASURE.ADJUSTMENT}

\subsection{편의 함수}

\begin{lstlisting}[language=Python]
from oecd_api import *

# GDP
df = get_gdp_quarterly("KOR+USA+JPN", "2020-Q1")     # 분기 GDP 성장률
df = get_gdp_annual("KOR+USA+JPN", "2020")            # 연간 GDP 성장률
df = get_gdp_level("KOR+USA+JPN", "2020-Q1")          # GDP 수준 (USD)

# 물가/고용/금리
df = get_cpi("KOR+USA", "2020-01")                    # CPI (지수+전년비+전월비)
df = get_cpi_yoy("KOR+USA", "2020-01")                # CPI 전년비만
df = get_unemployment("KOR+USA", "2020-01")            # 실업률
df = get_interest_rates("KOR+USA", "2020-01")          # 단기/장기 금리

# 기타
df = get_cli("KOR+USA", "2020-01")                    # 경기선행지수
df = get_industrial_production("KOR+USA", "2020-01")   # 산업생산지수
df = get_retail_sales("KOR+USA", "2020-01")            # 소매판매
df = get_trade("KOR+USA", "2020-01")                   # 수출/수입
\end{lstlisting}

\subsection{다운로드 항목 (8개)}

\begin{longtable}{l l}
\toprule
\textbf{파일명} & \textbf{설명} \\
\midrule
\endhead
\file{oecd\_gdp\_quarterly\_growth} & 분기 실질 GDP 성장률 (주요국) \\
\file{oecd\_gdp\_annual\_growth}    & 연간 GDP 성장률 (주요국) \\
\file{oecd\_gdp\_level\_usd}        & 분기 GDP 수준 (USD, 주요국) \\
\file{oecd\_cpi}                    & CPI 전년비 + 지수 (주요국) \\
\file{oecd\_unemployment}           & 조화실업률 (주요국, 월별) \\
\file{oecd\_interest\_rates}        & 단기/장기 금리 (주요국, 월별) \\
\file{oecd\_cli}                    & 경기선행지수 (주요국, 월별) \\
\file{oecd\_industrial\_production} & 산업생산지수 (주요국, 월별) \\
\bottomrule
\end{longtable}

\subsection{주의사항}

\begin{itemize}[leftmargin=2em]
    \item Rate Limit: 분당 약 10회 제한, 429 응답 시 자동 재시도 (10/20/30초 대기)
    \item \code{download\_global.py}는 항목 간 1초 딜레이 적용
    \item 중국(\code{CHN})은 일부 데이터플로우에서 누락될 수 있음
\end{itemize}


% ══════════════════════════════════════════════════════════════
\section{BIS API (\file{bis\_api.py})}
% ══════════════════════════════════════════════════════════════

\subsection{개요}

\textbf{BIS} (Bank for International Settlements, 국제결제은행)는 SDMX v2 REST API를 통해
글로벌 금융·거시 데이터를 제공합니다. 인증 불필요.

\begin{itemize}[leftmargin=2em]
    \item \textbf{기본 URL}: \code{https://stats.bis.org/api/v2}
    \item \textbf{응답 형식}: CSV (\code{format=csv})
    \item \textbf{국가 코드}: ISO 2자리 (KR, US, JP, DE, GB, FR, CN, ...)
\end{itemize}

\subsection{URL 구조}

\apibox{https://stats.bis.org/api/v2/data/dataflow/BIS/\{dataflow\}/\{version\}/\{key\}?format=csv\&startPeriod=YYYY-MM}

\subsection{데이터플로우 목록}

\begin{table}[h]
\centering
\caption{BIS 주요 데이터플로우}
\begin{tabularx}{\textwidth}{l l l X}
\toprule
\textbf{이름} & \textbf{Dataflow ID} & \textbf{버전} & \textbf{설명} \\
\midrule
CPI           & WS\_LONG\_CPI    & 1.0 & 소비자물가지수 \\
정책금리       & WS\_CBPOL        & 1.0 & 중앙은행 정책금리 \\
총 신용        & WS\_TC           & 2.0 & 비금융부문 총 신용 \\
신용/GDP 갭    & WS\_CREDIT\_GAP  & 1.0 & 바젤III 경기대응완충자본 지표 \\
부동산가격     & WS\_SPP          & 1.0 & 주택가격지수 \\
실효환율       & WS\_EER          & 1.0 & 명목/실질 실효환율 \\
DSR            & WS\_DSR          & 1.0 & 원리금상환비율 \\
\bottomrule
\end{tabularx}
\end{table}

\subsection{CPI 단위 코드}

\begin{table}[h]
\centering
\begin{tabular}{l l}
\toprule
\textbf{UNIT\_MEASURE} & \textbf{의미} \\
\midrule
\code{628} & Index (2010=100) \\
\code{771} & Year-on-year changes (\%) \\
\bottomrule
\end{tabular}
\end{table}

\subsection{신용 데이터 차입자 구분}

\begin{table}[h]
\centering
\begin{tabular}{l l}
\toprule
\textbf{Borrower 코드} & \textbf{의미} \\
\midrule
\code{P} & 민간 비금융부문 전체 \\
\code{H} & 가계 (Households) \\
\code{N} & 비금융기업 (Non-financial corporates) \\
\code{G} & 정부 (General government) \\
\code{C} & 핵심 부채 (P+G) \\
\bottomrule
\end{tabular}
\end{table}

\subsection{편의 함수}

\begin{lstlisting}[language=Python]
from bis_api import *

# 물가
df = get_cpi_yoy("KR+US+JP", "2020-01")      # CPI 전년비 (%)
df = get_cpi_index("KR+US+JP", "2020-01")     # CPI 지수 (2010=100)

# 금리
df = get_policy_rate("KR+US+JP", "2020-01")   # 중앙은행 정책금리

# 신용
df = get_credit_to_gdp("KR+US", "2015-Q1")                  # 민간 신용/GDP
df = get_credit_to_gdp("KR+US", "2015-Q1", borrower="H")    # 가계 신용/GDP
df = get_credit_to_gdp("KR+US", "2015-Q1", borrower="N")    # 기업 신용/GDP
df = get_credit_gap("KR+US", "2015-Q1")                      # 신용/GDP 갭

# 부동산/환율/DSR
df = get_property_prices("KR+US", "2015-Q1")                 # 주택가격지수
df = get_effective_exchange_rate("KR+US", "2020-01")          # 실효환율
df = get_debt_service_ratio("KR+US", "2015-Q1")              # 원리금상환비율
\end{lstlisting}

\subsection{다운로드 항목 (12개)}

\begin{longtable}{l l}
\toprule
\textbf{파일명} & \textbf{설명} \\
\midrule
\endhead
\file{bis\_cpi\_yoy}              & CPI 전년비 (\%, 주요국, 월별) \\
\file{bis\_cpi\_index}            & CPI 지수 (2010=100, 주요국, 월별) \\
\file{bis\_policy\_rate}          & 중앙은행 정책금리 (주요국, 월별) \\
\file{bis\_credit\_private\_gdp}  & 민간 신용/GDP (\%, 분기) \\
\file{bis\_credit\_household\_gdp}& 가계 신용/GDP (\%, 분기) \\
\file{bis\_credit\_corporate\_gdp}& 기업 신용/GDP (\%, 분기) \\
\file{bis\_credit\_govt\_gdp}     & 정부 신용/GDP (\%, 분기) \\
\file{bis\_credit\_gap}           & 신용/GDP 갭 (분기) \\
\file{bis\_property\_prices}      & 주택가격지수 (명목+실질, 분기) \\
\file{bis\_eer}                   & 실효환율 (명목+실질, 월별) \\
\file{bis\_dsr}                   & 원리금상환비율 (분기) \\
\bottomrule
\end{longtable}


% ══════════════════════════════════════════════════════════════
\section{IMF DataMapper API (\file{imf\_api.py})}
% ══════════════════════════════════════════════════════════════

\subsection{개요}

\textbf{IMF DataMapper API}는 World Economic Outlook (WEO) 데이터를 제공합니다.
연간 데이터 중심이며, 전망치(projection)도 포함됩니다. 인증 불필요.

\begin{itemize}[leftmargin=2em]
    \item \textbf{기본 URL}: \code{https://www.imf.org/external/datamapper/api/v1}
    \item \textbf{응답 형식}: JSON
    \item \textbf{국가 코드}: ISO 3자리 (ISO 2자리 자동 변환 지원)
\end{itemize}

\subsection{URL 구조}

\apibox{https://www.imf.org/external/datamapper/api/v1/\{INDICATOR\}/\{COUNTRIES\}}

응답 구조:
\begin{lstlisting}
{"values": {"INDICATOR": {"COUNTRY": {"YEAR": value, ...}, ...}}}
\end{lstlisting}

\subsection{주요 지표 코드}

\begin{table}[h]
\centering
\caption{IMF WEO 지표}
\begin{tabularx}{\textwidth}{l l X}
\toprule
\textbf{코드} & \textbf{분류} & \textbf{설명} \\
\midrule
\code{NGDP\_RPCH}    & 성장 & 실질 GDP 성장률 (\%) \\
\code{NGDPD}         & 성장 & 명목 GDP (십억 USD) \\
\code{NGDPDPC}       & 성장 & 1인당 GDP (USD) \\
\code{PPPGDP}        & 성장 & GDP PPP (십억 국제달러) \\
\code{PPPPC}         & 성장 & 1인당 GDP PPP (국제달러) \\
\code{PCPIPCH}       & 물가 & CPI 인플레이션 (\%, 연평균) \\
\code{PCPIEPCH}      & 물가 & CPI 인플레이션 (\%, 기말) \\
\code{LUR}           & 고용 & 실업률 (\%) \\
\code{LP}            & 고용 & 인구 (백만명) \\
\code{GGXWDG\_NGDP}  & 재정 & 정부 총부채/GDP (\%) \\
\code{GGXCNL\_NGDP}  & 재정 & 재정수지/GDP (\%) \\
\code{BCA\_NGDPD}    & 대외 & 경상수지/GDP (\%) \\
\bottomrule
\end{tabularx}
\end{table}

\subsection{지역 그룹 코드}

\begin{table}[h]
\centering
\begin{tabular}{l l}
\toprule
\textbf{코드} & \textbf{설명} \\
\midrule
\code{ADVEC}    & 선진국 \\
\code{OEMDC}    & 신흥시장·개발도상국 \\
\code{WEOWORLD} & 세계 전체 \\
\bottomrule
\end{tabular}
\end{table}

\subsection{편의 함수}

\begin{lstlisting}[language=Python]
from imf_api import *

# 성장
df = get_gdp_growth("USA/KOR/JPN", start_year=2010)
df = get_gdp_nominal("USA/KOR/JPN", start_year=2010)
df = get_gdp_per_capita("US/KR/JP", start_year=2010)  # ISO 2자리도 가능
df = get_gdp_ppp("USA/KOR", start_year=2010)

# 물가/고용
df = get_inflation("USA/KOR/JPN", start_year=2010)
df = get_unemployment("USA/KOR/JPN", start_year=2010)
df = get_population("USA/KOR/JPN", start_year=2010)

# 재정/대외
df = get_govt_debt("USA/KOR/JPN", start_year=2010)
df = get_fiscal_balance("USA/KOR/JPN", start_year=2010)
df = get_current_account("USA/KOR/JPN", start_year=2010)

# 범용
df = get_indicator("PPPPC", "USA/KOR", start_year=2010)

# 그룹
df = get_group_data("NGDP_RPCH", "WEOWORLD")         # 세계 GDP 성장률
df = get_country_snapshot("KOR", start_year=2020)     # 국가 종합
\end{lstlisting}

\subsection{다운로드 항목 (12개)}

\begin{longtable}{l l}
\toprule
\textbf{파일명} & \textbf{설명} \\
\midrule
\endhead
\file{imf\_gdp\_growth}           & 실질 GDP 성장률 (\%, 연간) \\
\file{imf\_gdp\_nominal}          & 명목 GDP (십억 USD, 연간) \\
\file{imf\_gdp\_per\_capita}      & 1인당 GDP (USD, 연간) \\
\file{imf\_gdp\_ppp}              & GDP PPP (십억 국제달러, 연간) \\
\file{imf\_inflation}             & CPI 인플레이션 (\%, 연간) \\
\file{imf\_unemployment}          & 실업률 (\%, 연간) \\
\file{imf\_population}            & 인구 (백만명, 연간) \\
\file{imf\_govt\_debt}            & 정부 총부채/GDP (\%, 연간) \\
\file{imf\_fiscal\_balance}       & 재정수지/GDP (\%, 연간) \\
\file{imf\_current\_account}      & 경상수지/GDP (\%, 연간) \\
\file{imf\_gdp\_per\_capita\_ppp} & 1인당 GDP PPP (국제달러) \\
\file{imf\_inflation\_eop}        & CPI 인플레이션 기말 (\%) \\
\bottomrule
\end{longtable}

\subsection{주의사항}

\begin{itemize}[leftmargin=2em]
    \item DataMapper API에서 사용 불가능한 지표: \code{GGR\_NGDP}, \code{GGX\_NGDP}, \code{TX\_RPCH}, \code{TM\_RPCH}
    \item WEO 데이터는 연 2회(4월/10월) 업데이트됨
    \item 미래 전망치(projection)도 포함되어 있으므로 실적치와 구분 필요
\end{itemize}


% ══════════════════════════════════════════════════════════════
\section{World Bank API (\file{worldbank\_api.py})}
% ══════════════════════════════════════════════════════════════

\subsection{개요}

\textbf{World Bank Open Data API v2}는 200+개국의 개발·거시경제 지표를 제공합니다.
인증 불필요.

\begin{itemize}[leftmargin=2em]
    \item \textbf{기본 URL}: \code{https://api.worldbank.org/v2}
    \item \textbf{응답 형식}: JSON (\code{format=json})
    \item \textbf{국가 코드}: ISO 2자리 (세미콜론 구분, 예: \code{KR;US;JP})
    \item \textbf{페이지네이션}: 자동 처리 (\code{per\_page=1000})
\end{itemize}

\subsection{URL 구조}

\apibox{https://api.worldbank.org/v2/country/\{countries\}/indicator/\{indicator\}?format=json\&per\_page=1000}

\subsection{주요 지표 코드}

\begin{longtable}{l l l}
\toprule
\textbf{분류} & \textbf{지표 코드} & \textbf{설명} \\
\midrule
\endhead
성장 & \code{NY.GDP.MKTP.CD}       & GDP (경상 US\$) \\
성장 & \code{NY.GDP.MKTP.KD.ZG}    & GDP 성장률 (연간 \%) \\
성장 & \code{NY.GDP.PCAP.CD}       & 1인당 GDP (경상 US\$) \\
성장 & \code{NY.GDP.MKTP.PP.CD}    & GDP PPP (경상 국제\$) \\
물가 & \code{FP.CPI.TOTL.ZG}       & 소비자물가 상승률 (연간 \%) \\
물가 & \code{NY.GDP.DEFL.KD.ZG}    & GDP 디플레이터 (연간 \%) \\
고용 & \code{SL.UEM.TOTL.ZS}       & 실업률 (ILO 기준, \%) \\
고용 & \code{SL.UEM.1524.ZS}       & 청년실업률 (15--24세, \%) \\
인구 & \code{SP.POP.TOTL}          & 총인구 \\
인구 & \code{SP.DYN.LE00.IN}       & 기대수명 (년) \\
무역 & \code{NE.TRD.GNFS.ZS}       & 무역/GDP (\%) \\
무역 & \code{BX.KLT.DINV.CD.WD}    & FDI 순유입 (US\$) \\
재정 & \code{GC.DOD.TOTL.GD.ZS}    & 정부부채/GDP (\%) \\
금융 & \code{FM.LBL.BMNY.GD.ZS}    & 광의통화/GDP (\%) \\
금융 & \code{FS.AST.PRVT.GD.ZS}    & 민간 국내신용/GDP (\%) \\
\bottomrule
\end{longtable}

\subsection{편의 함수}

\begin{lstlisting}[language=Python]
from worldbank_api import *

# 편의 함수
df = get_gdp_growth("KR;US;JP", start_year=2010)
df = get_inflation("KR;US;JP", start_year=2010)
df = get_unemployment("KR;US;JP", start_year=2010)
df = get_population("KR;US;JP", start_year=2010)

# 범용 함수
df = get_indicator("NY.GDP.MKTP.CD", "KR;US;JP", 2010, 2025)

# Wide format
wide = get_indicator_wide("FP.CPI.TOTL.ZG", "KR;US;JP", 2010)

# 국가 스냅샷 (주요 9개 지표)
df = get_country_snapshot("KR", start_year=2015)

# 지표 검색
df = search_indicators("GDP")

# 국가 목록
df = list_countries()
\end{lstlisting}

\subsection{다운로드 항목 (11개)}

\begin{longtable}{l l}
\toprule
\textbf{파일명} & \textbf{설명} \\
\midrule
\endhead
\file{wb\_gdp\_growth}             & GDP 성장률 (\%, 연간) \\
\file{wb\_gdp\_current\_usd}       & 명목 GDP (current USD, 연간) \\
\file{wb\_gdp\_per\_capita}        & 1인당 GDP (current USD, 연간) \\
\file{wb\_gdp\_ppp}                & GDP PPP (current intl\$, 연간) \\
\file{wb\_inflation}               & CPI 인플레이션 (\%, 연간) \\
\file{wb\_unemployment}            & 실업률 (ILO 추정, \%, 연간) \\
\file{wb\_population}              & 인구 (연간) \\
\file{wb\_trade\_pct\_gdp}         & 무역/GDP (\%, 연간) \\
\file{wb\_fdi\_net\_inflows}       & FDI 순유입 (current USD, 연간) \\
\file{wb\_broad\_money\_pct\_gdp}  & 광의통화/GDP (\%, 연간) \\
\file{wb\_domestic\_credit\_pct\_gdp} & 민간 국내신용/GDP (\%, 연간) \\
\bottomrule
\end{longtable}


% ══════════════════════════════════════════════════════════════
\section{CIA Factbook (\file{factbook\_api.py})}
% ══════════════════════════════════════════════════════════════

\subsection{개요}

\textbf{CIA World Factbook}은 261개 국가·지역의 경제 프로필을 JSON으로 제공합니다.
GitHub 리포지토리에서 매주 목요일 자동 업데이트됩니다.

\begin{itemize}[leftmargin=2em]
    \item \textbf{소스}: \url{https://github.com/factbook/factbook.json}
    \item \textbf{기본 URL}: \code{https://raw.githubusercontent.com/factbook/factbook.json/master}
    \item \textbf{국가 코드}: GEC(FIPS) 코드 --- ISO와 다름
    \item \textbf{데이터 유형}: 텍스트 기반 스냅샷 (최근 2--3년)
\end{itemize}

\subsection{GEC 국가 코드 매핑}

\begin{table}[h]
\centering
\caption{주요 국가 GEC 코드}
\begin{tabular}{l l l | l l l}
\toprule
\textbf{GEC} & \textbf{ISO} & \textbf{국가} & \textbf{GEC} & \textbf{ISO} & \textbf{국가} \\
\midrule
us & US & 미국     & fr & FR & 프랑스 \\
ks & KR & 한국     & ch & CN & 중국 \\
ja & JP & 일본     & it & IT & 이탈리아 \\
gm & DE & 독일     & ca & CA & 캐나다 \\
uk & GB & 영국     & as & AU & 호주 \\
mx & MX & 멕시코   & br & BR & 브라질 \\
in & IN & 인도     & id & ID & 인도네시아 \\
tu & TR & 튀르키예 & rs & RU & 러시아 \\
\bottomrule
\end{tabular}
\end{table}

\textbf{참고}: ISO 2자리 코드도 자동 변환됩니다 (\code{get\_economy\_summary("KR")} 가능).

\subsection{텍스트 파싱}

Factbook 데이터는 텍스트 형태로 제공되어 \code{\_parse\_number()} 함수로 숫자를 추출합니다:

\begin{lstlisting}[language=Python]
# 내부 파싱 예시
_parse_number("$25.676 trillion (2024 est.)")    # -> 25676000000000
_parse_number("2.8% (2024 est.)")                # -> 2.8
_parse_number("338,016,259 (2025 est.)")         # -> 338016259
_parse_number("52.3% of GDP (2023 est.)")        # -> 52.3
\end{lstlisting}

\subsection{편의 함수}

\begin{lstlisting}[language=Python]
from factbook_api import *

# 단일 국가 경제 요약 (dict)
summary = get_economy_summary("ks")     # GEC 코드
summary = get_economy_summary("KR")     # ISO 코드도 가능
# -> {"country": "ks", "gdp_ppp": ..., "gdp_growth": ...,
#     "inflation": ..., "unemployment": ..., ...}

# 여러 국가 비교 (DataFrame)
df = get_multi_country_summary(["us", "ks", "ja", "gm", "ch"])
df = get_multi_country_summary()     # 기본 15개국

# 시계열 (최근 2-3년)
df = get_gdp_growth_series("ks")
df = get_inflation_series("ks")
df = get_unemployment_series("ks")

# Raw 데이터
raw = get_country_raw("ks")          # 전체 JSON
econ = get_economy("ks")             # Economy 섹션만
\end{lstlisting}

\subsection{다운로드 항목 (1개)}

\begin{longtable}{l l}
\toprule
\textbf{파일명} & \textbf{설명} \\
\midrule
\file{factbook\_economy\_summary} & 주요국 경제 요약 (15개국 최신 스냅샷) \\
\bottomrule
\end{longtable}

\subsection{제공 필드}

\code{get\_economy\_summary()} 반환 dict의 주요 키:

\begin{table}[h]
\centering
\begin{tabularx}{\textwidth}{l X}
\toprule
\textbf{키} & \textbf{설명} \\
\midrule
\code{gdp\_ppp}                & 실질 GDP PPP (USD) \\
\code{gdp\_growth}             & GDP 성장률 (\%) \\
\code{gdp\_per\_capita}        & 1인당 GDP (USD) \\
\code{gdp\_official}           & GDP 명목 (USD) \\
\code{inflation}               & 인플레이션 (\%) \\
\code{unemployment}            & 실업률 (\%) \\
\code{population}              & 인구 \\
\code{labor\_force}            & 노동력 \\
\code{public\_debt\_pct\_gdp}  & 정부부채/GDP (\%) \\
\code{current\_account}        & 경상수지 (USD) \\
\code{exports} / \code{imports}& 수출 / 수입 (USD) \\
\code{fx\_reserves}            & 외환보유고 (USD) \\
\code{gdp\_share\_agriculture} & 농업 비중 (\%) \\
\code{gdp\_share\_industry}    & 산업 비중 (\%) \\
\code{gdp\_share\_services}    & 서비스 비중 (\%) \\
\bottomrule
\end{tabularx}
\end{table}


% ══════════════════════════════════════════════════════════════
\section{한국은행 ECOS API (\file{ecos\_api.py})}
% ══════════════════════════════════════════════════════════════

\subsection{개요}

\textbf{한국은행 ECOS Open API}는 한국의 경제통계를 제공합니다.
별도 다운로드 스크립트 (\file{download\_major.py}, \file{download\_extra.py}, \file{download\_remaining.py})를
통해 163개 통계 테이블의 parquet 파일이 \file{data/}에 저장됩니다.

\begin{itemize}[leftmargin=2em]
    \item \textbf{공식 사이트}: \url{https://ecos.bok.or.kr/api/}
    \item \textbf{기본 URL}: \code{https://ecos.bok.or.kr/api}
    \item \textbf{API 키}: 코드에 내장 (하드코딩)
\end{itemize}

\subsection{URL 구조}

\apibox{https://ecos.bok.or.kr/api/\{service\}/\{api\_key\}/json/kr/1/100000/\{params\}}

\subsection{주요 서비스}

\begin{table}[h]
\centering
\begin{tabular}{l l}
\toprule
\textbf{서비스} & \textbf{설명} \\
\midrule
\code{StatisticTableList} & 통계표 목록 조회 \\
\code{StatisticItemList}  & 통계 세부항목 목록 \\
\code{StatisticSearch}    & 통계 데이터 조회 \\
\bottomrule
\end{tabular}
\end{table}

\subsection{ECOS 다운로드}

ECOS 데이터는 \code{download\_global.py}와 별도로 운영됩니다:

\begin{lstlisting}[style=bash]
# ECOS 전용 다운로드 스크립트
python3 download_major.py        # 주요 통계
python3 download_extra.py        # 추가 통계
python3 download_remaining.py    # 나머지 통계

# 저장 위치: data/ (163개 parquet 파일)
\end{lstlisting}


% ══════════════════════════════════════════════════════════════
\section{데이터 활용 예시}
% ══════════════════════════════════════════════════════════════

\subsection{파일 로드}

\begin{lstlisting}[language=Python]
import pandas as pd

# Parquet 로드 (pyarrow 필요, 빠르고 타입 보존)
df = pd.read_parquet("data/fred_gdp_real.parquet")
print(df.columns)  # ['date', 'value', 'series_id']
print(df.tail())

# CSV 로드 (추가 패키지 불필요, Excel 호환)
df = pd.read_csv("data/fred_gdp_real.csv")
print(df.tail())

# 여러 파일 병합
import glob
files = glob.glob("data/imf_*.parquet")
for f in files:
    df = pd.read_parquet(f)
    print(f"{f}: {len(df)} rows, {list(df.columns)}")
\end{lstlisting}

\textbf{Parquet vs CSV 비교:}
\begin{table}[h]
\centering
\begin{tabularx}{\textwidth}{l X X}
\toprule
& \textbf{Parquet} & \textbf{CSV} \\
\midrule
용량     & 작음 (컬럼 압축)         & 큼 (텍스트 원본) \\
속도     & 빠름                     & 느림 \\
타입 보존 & 날짜·숫자 타입 유지       & 문자열로 저장 (파싱 필요) \\
호환성   & Python/R/Spark           & Excel/텍스트 편집기/모든 도구 \\
의존성   & pyarrow 또는 fastparquet & 없음 \\
\bottomrule
\end{tabularx}
\end{table}

\subsection{국가간 GDP 성장률 비교}

\begin{lstlisting}[language=Python]
import pandas as pd

# OECD 분기 GDP 성장률
df = pd.read_parquet("data/oecd_gdp_quarterly_growth.parquet")

# 한국 vs 미국 vs 일본 필터
mask = df["REF_AREA"].isin(["KOR", "USA", "JPN"])
subset = df.loc[mask, ["REF_AREA", "TIME_PERIOD", "OBS_VALUE"]]

# Wide format으로 변환
wide = subset.pivot(index="TIME_PERIOD", columns="REF_AREA", values="OBS_VALUE")
print(wide.tail(8))
\end{lstlisting}

\subsection{BIS 신용 데이터 분석}

\begin{lstlisting}[language=Python]
import pandas as pd

# 민간 신용/GDP
private = pd.read_parquet("data/bis_credit_private_gdp.parquet")

# 가계 신용/GDP
household = pd.read_parquet("data/bis_credit_household_gdp.parquet")

# 한국 가계부채 추이
kr = household[household["BORROWERS_CTY"] == "KR"]
print(kr[["TIME_PERIOD", "OBS_VALUE"]].tail(10))
\end{lstlisting}

\subsection{다중 소스 교차 검증}

\begin{lstlisting}[language=Python]
import pandas as pd

# GDP 성장률: IMF vs World Bank 비교
imf = pd.read_parquet("data/imf_gdp_growth.parquet")
wb = pd.read_parquet("data/wb_gdp_growth.parquet")

# 한국 2020년 데이터 비교
imf_kr = imf[(imf["country"] == "KOR") & (imf["year"] == 2020)]
wb_kr = wb[(wb["country_id"] == "KR") & (wb["date"] == 2020)]

print(f"IMF: {imf_kr['value'].values}")
print(f"WB:  {wb_kr['value'].values}")
\end{lstlisting}


% ══════════════════════════════════════════════════════════════
\section{전체 다운로드 파일 목록}
% ══════════════════════════════════════════════════════════════

\code{download\_global.py}가 생성하는 전체 parquet 파일 목록입니다 (총 61개).

\begin{longtable}{r l l l}
\toprule
\textbf{\#} & \textbf{파일명} & \textbf{소스} & \textbf{주기} \\
\midrule
\endhead
\multicolumn{4}{l}{\textbf{FRED (17개)}} \\
\midrule
1  & fred\_gdp\_real              & FRED & 분기 \\
2  & fred\_gdp\_growth            & FRED & 분기 \\
3  & fred\_cpi                    & FRED & 월별 \\
4  & fred\_cpi\_core              & FRED & 월별 \\
5  & fred\_pce                    & FRED & 월별 \\
6  & fred\_pce\_core              & FRED & 월별 \\
7  & fred\_unemployment           & FRED & 월별 \\
8  & fred\_payrolls               & FRED & 월별 \\
9  & fred\_participation          & FRED & 월별 \\
10 & fred\_initial\_claims        & FRED & 주별 \\
11 & fred\_fed\_funds             & FRED & 월별 \\
12 & fred\_treasury\_10y          & FRED & 월별 \\
13 & fred\_treasury\_2y           & FRED & 월별 \\
14 & fred\_term\_spread           & FRED & 월별 \\
15 & fred\_industrial\_prod       & FRED & 월별 \\
16 & fred\_retail\_sales          & FRED & 월별 \\
17 & fred\_consumer\_sentiment    & FRED & 월별 \\
18 & fred\_m2                     & FRED & 월별 \\
\midrule
\multicolumn{4}{l}{\textbf{OECD (8개)}} \\
\midrule
19 & oecd\_gdp\_quarterly\_growth & OECD & 분기 \\
20 & oecd\_gdp\_annual\_growth    & OECD & 연간 \\
21 & oecd\_gdp\_level\_usd        & OECD & 분기 \\
22 & oecd\_cpi                    & OECD & 월별 \\
23 & oecd\_unemployment           & OECD & 월별 \\
24 & oecd\_interest\_rates        & OECD & 월별 \\
25 & oecd\_cli                    & OECD & 월별 \\
26 & oecd\_industrial\_production & OECD & 월별 \\
\midrule
\multicolumn{4}{l}{\textbf{BIS (11개)}} \\
\midrule
27 & bis\_cpi\_yoy                & BIS & 월별 \\
28 & bis\_cpi\_index              & BIS & 월별 \\
29 & bis\_policy\_rate            & BIS & 월별 \\
30 & bis\_credit\_private\_gdp    & BIS & 분기 \\
31 & bis\_credit\_household\_gdp  & BIS & 분기 \\
32 & bis\_credit\_corporate\_gdp  & BIS & 분기 \\
33 & bis\_credit\_govt\_gdp       & BIS & 분기 \\
34 & bis\_credit\_gap             & BIS & 분기 \\
35 & bis\_property\_prices        & BIS & 분기 \\
36 & bis\_eer                     & BIS & 월별 \\
37 & bis\_dsr                     & BIS & 분기 \\
\midrule
\multicolumn{4}{l}{\textbf{IMF (12개)}} \\
\midrule
38 & imf\_gdp\_growth             & IMF & 연간 \\
39 & imf\_gdp\_nominal            & IMF & 연간 \\
40 & imf\_gdp\_per\_capita        & IMF & 연간 \\
41 & imf\_gdp\_ppp                & IMF & 연간 \\
42 & imf\_inflation               & IMF & 연간 \\
43 & imf\_unemployment            & IMF & 연간 \\
44 & imf\_population              & IMF & 연간 \\
45 & imf\_govt\_debt              & IMF & 연간 \\
46 & imf\_fiscal\_balance         & IMF & 연간 \\
47 & imf\_current\_account        & IMF & 연간 \\
48 & imf\_gdp\_per\_capita\_ppp   & IMF & 연간 \\
49 & imf\_inflation\_eop          & IMF & 연간 \\
\midrule
\multicolumn{4}{l}{\textbf{World Bank (11개)}} \\
\midrule
50 & wb\_gdp\_growth              & WB & 연간 \\
51 & wb\_gdp\_current\_usd        & WB & 연간 \\
52 & wb\_gdp\_per\_capita         & WB & 연간 \\
53 & wb\_gdp\_ppp                 & WB & 연간 \\
54 & wb\_inflation                & WB & 연간 \\
55 & wb\_unemployment             & WB & 연간 \\
56 & wb\_population               & WB & 연간 \\
57 & wb\_trade\_pct\_gdp          & WB & 연간 \\
58 & wb\_fdi\_net\_inflows        & WB & 연간 \\
59 & wb\_broad\_money\_pct\_gdp   & WB & 연간 \\
60 & wb\_domestic\_credit\_pct\_gdp & WB & 연간 \\
\midrule
\multicolumn{4}{l}{\textbf{CIA Factbook (1개)}} \\
\midrule
61 & factbook\_economy\_summary   & CIA & 스냅샷 \\
\bottomrule
\end{longtable}


% ══════════════════════════════════════════════════════════════
\section{트러블슈팅}
% ══════════════════════════════════════════════════════════════

\subsection{공통 문제}

\begin{description}[leftmargin=2em]
    \item[FRED API 키 오류] \code{FRED\_API\_KEY} 환경변수를 설정하거나, 키가 유효한지 확인하세요.
        키 없이 실행하면 FRED는 자동으로 건너뜁니다.

    \item[OECD 429 Rate Limit] OECD는 분당 요청 수를 제한합니다.
        스크립트가 자동으로 10/20/30초 대기 후 재시도합니다. 연속 실행 시 간격을 두세요.

    \item[빈 DataFrame 반환] 특정 국가·기간에 데이터가 없을 수 있습니다.
        스크립트는 빈 결과를 ``데이터 없음''으로 표시하고 건너뜁니다.

    \item[네트워크 타임아웃] 각 API 모듈은 30--60초 타임아웃을 설정합니다.
        네트워크가 불안정하면 개별 소스를 선택적으로 재실행하세요.

    \item[Parquet 의존성] \code{pyarrow} 또는 \code{fastparquet}가 필요합니다:
\end{description}

\begin{lstlisting}[style=bash]
pip install pyarrow
# 또는
pip install fastparquet
\end{lstlisting}

\subsection{소스별 특이사항}

\begin{table}[h]
\centering
\begin{tabularx}{\textwidth}{l X}
\toprule
\textbf{소스} & \textbf{주의사항} \\
\midrule
OECD & CHN 등 비회원국은 일부 데이터플로우에서 누락 가능 \\
BIS  & 일부 국가(신흥국)는 데이터 제공 기간이 짧을 수 있음 \\
IMF  & WEO 전망치가 포함됨 --- 실적치와 구분 필요 \\
WB   & 최신 연도 데이터가 1--2년 지연될 수 있음 \\
Factbook & 텍스트 파싱 기반이므로 형식 변경 시 오류 가능 \\
ECOS & \code{download\_global.py}와 별도 운영 (\code{download\_major/extra/remaining.py}) \\
\bottomrule
\end{tabularx}
\end{table}


% ══════════════════════════════════════════════════════════════
\section{부록: 국가 코드 대조표}
% ══════════════════════════════════════════════════════════════

\begin{longtable}{l l l l l}
\toprule
\textbf{국가} & \textbf{ISO 3} & \textbf{ISO 2} & \textbf{GEC(FIPS)} & \textbf{비고} \\
\midrule
\endhead
미국       & USA & US & us & \\
한국       & KOR & KR & ks & GEC 주의 \\
일본       & JPN & JP & ja & GEC 주의 \\
독일       & DEU & DE & gm & GEC 주의 \\
영국       & GBR & GB & uk & \\
프랑스     & FRA & FR & fr & \\
중국       & CHN & CN & ch & GEC 주의 \\
이탈리아   & ITA & IT & it & \\
캐나다     & CAN & CA & ca & \\
호주       & AUS & AU & as & GEC 주의 \\
멕시코     & MEX & MX & mx & \\
브라질     & BRA & BR & br & \\
인도       & IND & IN & in & \\
인도네시아 & IDN & ID & id & \\
튀르키예   & TUR & TR & tu & GEC 주의 \\
러시아     & RUS & RU & rs & GEC 주의 \\
남아프리카 & ZAF & ZA & za & \\
사우디     & SAU & SA & sa & \\
아르헨티나 & ARG & AR & ar & \\
나이지리아 & NGA & NG & ni & GEC 주의 \\
\bottomrule
\end{longtable}


\end{document}
